%------------------------------------------------------------------------------%
\documentclass[a4paper,12pt,fleqn]{article}

\usepackage{gaal}

\ShowAnswers

%------------------------------------------------------------------------------%
\begin{document}

\makeHeader{GAAL}{Prova 3 -- S}{2026-2}

% 01
%------------------------------------------------------------------------------%
\exercise{25\%}
Determine um vetor normal e uma equação cartesiana do plano que passa pelos pontos
\[
  P = (1,2,-1)
  \qquad
  Q = (3,1,0)
  \qquad
  R = (0,0,2)
\]

\begin{answer}
  Calculamos dois vetores contidos no plano

  \begin{align*}
    \overrightarrow{PQ}
    & = Q-P
    \\
    & = (3,1,0)-(1,2,-1)
    \\
    & = (2,-1,1)
  \end{align*}

  \begin{align*}
    \overrightarrow{PR}
    & = R-P
    \\
    & = (0,0,2)-(1,2,-1)
    \\
    & = (-1,-2,3)
  \end{align*}

  Um vetor normal ao plano é dado pelo produto vetorial

  \begin{align*}
    \vec{n}
    & = \overrightarrow{PQ}\times\overrightarrow{PR}
    \\
    & = \begin{vmatrix}
          \vec{i} & \vec{j} & \vec{k} \\
          2       & -1      & 1       \\
          -1      & -2      & 3
        \end{vmatrix}
    \\
    & = \vec{i}(-3+2)
      -\vec{j}(6+1)
      +\vec{k}(-4-1)
    \\
    & = -\vec{i}-7\vec{j}-5\vec{k}
    \\
    & = (-1,-7,-5)
  \end{align*}

  Podemos utilizar também o vetor normal oposto
  \[
    \vec{n}=(1,7,5)
  \]

  A equação do plano que passa por $P=(1,2,-1)$ é

  \begin{align*}
    \vec{n}\cdot(X-P)&=0
    \\
    (1,7,5)\cdot(x-1,y-2,z+1)&=0
    \\
    (x-1)+7(y-2)+5(z+1)&=0
    \\
    x-1+7y-14+5z+5&=0
    \\
    x+7y+5z&=10
  \end{align*}
\end{answer}

% 02
%------------------------------------------------------------------------------%
\exercise{25\%}
Determine a posição relativa entre os planos e obtenha uma equação paramétrica
da interseção
\[
  \pi_1\colon 2x-y+z=3
  \qquad
  \pi_2\colon x+y-2z=0
\]

\begin{answer}
  Os vetores normais dos planos são

  \[
    \vec{n}_1=(2,-1,1)
    \qquad
    \vec{n}_2=(1,1,-2)
  \]

  Como os vetores normais não são proporcionais, os planos são secantes
  e sua interseção é uma reta

  Um vetor diretor da reta é dado pelo produto vetorial dos vetores normais

  \begin{align*}
    \vec{d}
    & = \vec{n}_1\times\vec{n}_2
    \\
    & = \begin{vmatrix}
          \vec{i} & \vec{j} & \vec{k} \\
          2       & -1      & 1       \\
          1       & 1       & -2
        \end{vmatrix}
    \\
    & = \vec{i}(2-1)
      -\vec{j}(-4-1)
      +\vec{k}(2+1)
    \\
    & = (1,5,3)
  \end{align*}

  Para encontrar um ponto da reta, fazemos $z=0$

  \begin{align*}
    2x-y&=3
    \\
    x+y&=0
  \end{align*}

  Somando as equações

  \begin{align*}
    3x&=3
    \\
    x&=1
    \\
    y&=-1
  \end{align*}

  Portanto, o ponto $P=(1,-1,0)$ pertence à reta de interseção

  Uma equação paramétrica da reta é

  \[
    (x,y,z)=(1,-1,0)+t(1,5,3)
  \]

  Ou, equivalentemente,

  \[
    \begin{cases}
      x=1+t\\
      y=-1+5t\\
      z=3t
    \end{cases}
    \qquad t\in\mathbb{R}
  \]
\end{answer}

%% 03
%%------------------------------------------------------------------------------%
%\exercise{20\%}
%Determine o ângulo entre os vetores
%\[
%  \vec{v}=(1,2,-2)
%  \qquad
%  \vec{w}=(2,-1,1)
%\]
%
%\begin{answer}
%  Calculamos inicialmente o produto escalar
%
%  \begin{align*}
%    \vec{v}\cdot\vec{w}
%    & = 1\cdot 2+2\cdot(-1)+(-2)\cdot 1
%    \\
%    & = 2-2-2
%    \\
%    & = -2
%  \end{align*}
%
%  As normas dos vetores são
%
%  \begin{align*}
%    \norm{\vec{v}}
%    & = \sqrt{1^2+2^2+(-2)^2}
%    \\
%    & = \sqrt{1+4+4}
%    \\
%    & = 3
%  \end{align*}
%
%  \begin{align*}
%    \norm{\vec{w}}
%    & = \sqrt{2^2+(-1)^2+1^2}
%    \\
%    & = \sqrt{4+1+1}
%    \\
%    & = \sqrt{6}
%  \end{align*}
%
%  Pela definição de produto escalar
%
%  \begin{align*}
%    \cos(\theta)
%    & = \frac{\vec{v}\cdot\vec{w}}
%             {\norm{\vec{v}}\norm{\vec{w}}}
%    \\
%    & = \frac{-2}{3\sqrt{6}}
%  \end{align*}
%
%  Portanto
%
%  \[
%    \theta=\arccos\left(-\frac{2}{3\sqrt{6}}\right)
%  \]
%
%  \[
%    {\color{gray}
%      \theta\approx\ang{105.79}
%    }
%  \]
%\end{answer}

% 04
%------------------------------------------------------------------------------%
\exercise{25\%}
Verifique se os vetores
\[
  \vec{v}_1=(1,2,0)
  \qquad
  \vec{v}_2=(0,1,1)
  \qquad
  \vec{w}=(1,4,2)
\]
são linearmente independentes

Caso sejam linearmente dependentes, determine uma relação de dependência
linear entre eles

\begin{answer}
  Procuramos escalares $a$ e $b$ tais que

  \[
    a\vec{v}_1+b\vec{v}_2=\vec{w}
  \]

  Substituindo os vetores

  \[
    a(1,2,0)+b(0,1,1)=(1,4,2)
  \]

  Comparando as coordenadas, obtemos o sistema

  \[
    \begin{cases}
      a=1\\
      2a+b=4\\
      b=2
    \end{cases}
  \]

  A solução é $a=1$ e $b=2$

  Portanto

  \[
    \vec{w}=\vec{v}_1+2\vec{v}_2
  \]

  Assim, os três vetores são linearmente dependentes

  Uma relação de dependência linear é

  \[
    \vec{v}_1+2\vec{v}_2-\vec{w}=\vec{0}
  \]

  Além disso, os vetores $\vec{v}_1$ e $\vec{v}_2$ são
  linearmente independentes, pois

  \[
    a\vec{v}_1+b\vec{v}_2=\vec{0}
  \]

  implica

  \[
    \begin{cases}
      a=0\\
      2a+b=0\\
      b=0
    \end{cases}
  \]

  Logo, $a=b=0$
\end{answer}

% 05
%------------------------------------------------------------------------------%
\exercise{20\%}
Obtenha, pelo processo de Gram--Schmidt, uma base ortonormal que gere
o mesmo espaço que os vetores
\[
  \vec{v}_1=(1,1,0)
  \qquad
  \vec{v}_2=(1,0,1)
  \qquad
  \vec{v}_3=(0,1,1)
\]

\begin{answer}
  Aplicamos o processo de Gram--Schmidt

  **Primeiro vetor**

  \begin{align*}
    \vec{u}_1
    & = \vec{v}_1
    \\
    & = (1,1,0)
  \end{align*}

  Sua norma é

  \begin{align*}
    \norm{\vec{u}_1}
    & = \sqrt{1^2+1^2}
    \\
    & = \sqrt{2}
  \end{align*}

  Portanto

  \[
    \vec{e}_1
    = \frac{\vec{u}_1}{\norm{\vec{u}_1}}
    = \frac{1}{\sqrt{2}}(1,1,0)
  \]

  **Segundo vetor**

  \begin{align*}
    \vec{u}_2
    & = \vec{v}_2
      -\frac{\vec{v}_2\cdot\vec{u}_1}
             {\norm{\vec{u}_1}^2}\vec{u}_1
    \\
    & = (1,0,1)
      -\frac{1}{2}(1,1,0)
    \\
    & = \left(\frac{1}{2},-\frac{1}{2},1\right)
  \end{align*}

  Sua norma é

  \begin{align*}
    \norm{\vec{u}_2}
    & = \sqrt{\frac{1}{4}+\frac{1}{4}+1}
    \\
    & = \sqrt{\frac{3}{2}}
    \\
    & = \frac{\sqrt{6}}{2}
  \end{align*}

  Assim

  \begin{align*}
    \vec{e}_2
    & = \frac{\vec{u}_2}{\norm{\vec{u}_2}}
    \\
    & = \frac{2}{\sqrt{6}}
      \left(\frac{1}{2},-\frac{1}{2},1\right)
    \\
    & = \frac{1}{\sqrt{6}}(1,-1,2)
  \end{align*}

  **Terceiro vetor**

  \begin{align*}
    \vec{u}_3
    & = \vec{v}_3
      -\frac{\vec{v}_3\cdot\vec{u}_1}
             {\norm{\vec{u}_1}^2}\vec{u}_1
      -\frac{\vec{v}_3\cdot\vec{u}_2}
             {\norm{\vec{u}_2}^2}\vec{u}_2
    \\
    & = (0,1,1)
      -\frac{1}{2}(1,1,0)
      -\frac{1}{3}
       \left(\frac{1}{2},-\frac{1}{2},1\right)
    \\
    & = \left(-\frac{2}{3},\frac{2}{3},\frac{2}{3}\right)
  \end{align*}

  Sua norma é

  \begin{align*}
    \norm{\vec{u}_3}
    & = \sqrt{
        \frac{4}{9}+\frac{4}{9}+\frac{4}{9}
      }
    \\
    & = \sqrt{\frac{4}{3}}
    \\
    & = \frac{2}{\sqrt{3}}
  \end{align*}

  Portanto

  \begin{align*}
    \vec{e}_3
    & = \frac{\vec{u}_3}{\norm{\vec{u}_3}}
    \\
    & = \frac{\sqrt{3}}{2}
      \left(-\frac{2}{3},\frac{2}{3},\frac{2}{3}\right)
    \\
    & = \frac{1}{\sqrt{3}}(-1,1,1)
  \end{align*}

  A base ortonormal procurada é

  \[
    \left\{
      \frac{1}{\sqrt{2}}(1,1,0)
      \qquad
      \frac{1}{\sqrt{6}}(1,-1,2)
      \qquad
      \frac{1}{\sqrt{3}}(-1,1,1)
    \right\}
  \]
\end{answer}

%------------------------------------------------------------------------------%
\end{document}
%------------------------------------------------------------------------------%
